% Appendix B: Derivations (methods author). PLAN.md section 1.4: the proofs the sidebars only cite. B.1 Taylor
% expansions and the moment conditions, B.2 the Lagrange formula, B.3 Richardson extrapolation, B.4 the error model
% and the best step. No figures, no listings; the typed numbers are exact constants, the generated cost factors of
% B.4 are macros. Revised at stage 6 per reviews/R1-decision.md (T8, library-17, evidence-21); the changes are listed
% in responses/R1-a2-derivations.md.
\chapter{Derivations}\label{ch:derivations}

This appendix proves what the sidebars of \cref{ch:stencils,ch:orders,ch:ridders} only cite: why the moment
conditions decide a rule's order and its error constant, why the Lagrange formula produces the weights, why
Richardson's tableau works and what its error estimate measures, and where the best step and its constant come
from. Throughout, $f$ has as many continuous derivatives as the expansions use, the nodes $n_k$ are distinct
integers, and the weights $w_k$ are the rational numbers of \cref{eq:stencils-rule}; with the library's integer
numerators $c_k = \den w_k$ every moment is multiplied by $\den$.

\section{Taylor expansions and the moment conditions}\label{sec:derivations-moments}

Taylor's theorem with the Lagrange remainder\cite[Sect.~1.1]{burden2016} gives, for each node, with
$x_k = x + n_k h$,
\[
  f(x_k) = \sum_{j=0}^{m-1} \frac{f^{(j)}(x)}{j!}\, n_k^j h^j + \frac{f^{(m)}(\xi_k)}{m!}\, n_k^m h^m,
  \qquad \xi_k \text{ between } x \text{ and } x_k .
\]
Multiply by $w_k$ and sum over the nodes. The powers of $h$ collect into the moments
$\mom{j} = \sum_k w_k n_k^j$ of \cref{eq:stencils-moments}:
\begin{equation}
  \sum_k w_k f(x_k) = \sum_{j=0}^{m-1} \frac{\mom{j}}{j!}\, h^j f^{(j)}(x) + \frac{h^m}{m!} \sum_k w_k n_k^m f^{(m)}(\xi_k).
  \label{eq:derivations-taylor}
\end{equation}
Suppose the moment conditions hold up to $m - 1 = \dorder + \aorder - 1$: $\mom{\dorder} = \dorder!$ and
$\mom{j} = 0$ for every other $j$ below $m$. Then the sum on the left is $h^{\dorder} f^{(\dorder)}(x)$ plus the
remainder, and dividing by $h^{\dorder}$ gives
\[
  \frac{1}{h^{\dorder}} \sum_k w_k f(x_k) = f^{(\dorder)}(x) + \frac{h^{\aorder}}{(\dorder + \aorder)!}
  \sum_k w_k n_k^{\dorder + \aorder} f^{(\dorder + \aorder)}(\xi_k).
\]
The remainder is $\order(h^{\aorder})$, which is the accuracy order, and since each $\xi_k$ tends to $x$ with
$h$, the remainder is $\mom{\dorder + \aorder}\, h^{\aorder} f^{(\dorder + \aorder)}(x) / (\dorder + \aorder)! +
\order(h^{\aorder + 1})$ for every rule. This asymptotic form is what the error model uses, with the error constant
\begin{equation}
  \errc = \frac{\abs{\mom{\dorder + \aorder}}}{(\dorder + \aorder)!}, \qquad
  \text{or } \frac{\abs{\mom{\dorder + \aorder}}}{(\dorder + \aorder)!\,\den} \text{ with the integer numerators,}
  \label{eq:derivations-errc}
\end{equation}
as \cref{ch:tables} tabulates it.

The sharper statement with a single intermediate point, as in \cref{eq:stencils-central-four} and
\cref{eq:orders-second}, needs one more argument, and which argument depends on the rule. If the products
$w_k n_k^{\dorder + \aorder}$ all have one sign, the intermediate-value theorem collapses the sum above to one
point: the remainder is $\mom{\dorder + \aorder}\, h^{\aorder} f^{(\dorder + \aorder)}(\xi) / (\dorder + \aorder)!$
for some $\xi$ in the span of the stencil. That is the case for \cref{eq:orders-second}, whose products are
$1, 0, 1$, but not for \cref{eq:stencils-central-four}, whose products $-\tfrac83, \tfrac23, \tfrac23, -\tfrac83$
have mixed signs, nor for any central rule for $f'$ above order 2, nor for some one-sided rules. For the first
derivative the single-point form follows instead from the interpolation error formula, as Burden et
al.\cite[Sect.~4.1]{burden2016} derive the five-point formulas: with $x$ itself among the nodes (added with the
weight 0 where the rule omits it, which changes nothing), $f(x + th) - p(t)$ equals $f^{(m)}(\xi(t)) \prod_j (t -
n_j)\, h^m / m!$, and differentiating once at $t = 0$, where the node product vanishes, leaves $f'(x) - p'(0)/h =
f^{(m)}(\xi)\, h^{m-1} \prod_{n_j \ne 0} (-n_j) / m!$ with one $\xi$. For \code{Central<4>} the product over the
nodes $-2, -1, 1, 2$ is $4$ and $m = 5$, which gives the $+ h^4 f^{(5)}(\xi) / 30$ of
\cref{eq:stencils-central-four}; the one-sided rules, which contain $x$ as a node, are covered by the same formula.
For the higher derivatives and mixed signs the book uses only the asymptotic form.

For a symmetric stencil, nodes in pairs $\pm n$ with $w_{-n} = (-1)^{\dorder} w_n$, every moment $\mom{j}$ with
$j - \dorder$ odd vanishes identically, so the expansion of a central rule contains only the powers
$h^{\aorder}, h^{\aorder + 2}, \ldots$: an even expansion in $h$. This is what \cref{sb:remainder} states for
the central difference and what the extrapolation of \cref{sec:derivations-richardson} relies on for every
order. It also explains why a central rule needs one node fewer than a one-sided rule of the same order: of the
$\dorder + \aorder$ conditions, those with $j - \dorder$ odd hold by symmetry and cost no weight.

Rounding enters \cref{eq:derivations-taylor} through the samples. If each $f(x_k)$ is off by at most
$\epsmach \abs{f(x_k)}$, the weighted sum is off by at most $\sum_k \abs{w_k}\, \epsmach \abs{f(x_k)}$, which
for samples of about the size of $\abs{f(x)}$ is $\ampl \epsmach \abs{f}$ with $\ampl = \sum_k \abs{w_k}$; the
division by $h^{\dorder}$ turns it into the rounding term $\ampl \epsmach \abs{f} / h^{\dorder}$ of
\cref{eq:orders-model}. Rounding of the points, $x_k$ off by up to $\epsmach \abs{x_k} / 2$, moves each sample by
about $\epsmach \abs{x f'} / 2$ and adds the third term named in \cref{sec:stencils-price}.

\section{The Lagrange formula for the weights}\label{sec:derivations-lagrange}

Let $p$ be the polynomial of degree below the number of nodes $m$ that interpolates the samples, written in
the scaled variable $t$ with $x + t h$: $p(t) = \sum_k f(x_k) L_k(t)$, where the Lagrange basis polynomial
$L_k(t) = \prod_{j \ne k} (t - n_j) / (n_k - n_j)$ is 1 at $t = n_k$ and 0 at every other node. Differentiating
$p$ $\dorder$ times at $t = 0$ and dividing by $h^{\dorder}$ gives an estimate of $f^{(\dorder)}(x)$ of the form
\cref{eq:stencils-rule} with
\begin{equation}
  w_k = L_k^{(\dorder)}(0) = \dorder!\,[t^{\dorder}]\, L_k(t),
  \label{eq:derivations-lagrange}
\end{equation}
the coefficient of $t^{\dorder}$ in $L_k$ times $\dorder!$, which is what the library's generator computes with
exact integer arithmetic.\cite[Sect.~4.1]{burden2016} These weights satisfy the moment conditions, for the
following reason. Interpolation on $m$ nodes reproduces every polynomial of degree below $m$: if $f(x + t h) = t^j$
with $j < m$, then $p(t) = t^j$ exactly, since both sides are polynomials of degree below $m$ that agree at $m$
points. For that $f$ the samples are $f(x_k) = n_k^j$, so $\sum_k w_k n_k^j = p^{(\dorder)}(0) =
(t^j)^{(\dorder)}(0)$, which is $\dorder!$ for $j = \dorder$ and 0 for every other $j < m$. These are exactly the
moment conditions up to $\mom{m-1}$, so the Lagrange weights are the unique solution of that linear system, and
the rule's order is at least $m - \dorder$. For the one-sided rules $m = \aorder + \dorder$ and the order is
exactly $\aorder$, since $\mom{m}$ does not vanish; for the central rules $m = \aorder + \dorder - 1$ and the
symmetry argument of \cref{sec:derivations-moments} supplies the missing condition $\mom{m} = 0$.

The library's own worked example is \code{Central<4>} for $f'$ on the nodes $-2, -1, 1, 2$. For the node
$n_k = -1$, the numerator of $L_k$ is $(t + 2)(t - 1)(t - 2) = t^3 - t^2 - 4t + 4$ and its denominator
$(-1 + 2)(-1 - 1)(-1 - 2) = 6$, so $w_{-1} = 1! \cdot (-4) / 6 = -2/3$; the other three nodes give $1/12$,
$2/3$ and $-1/12$. Over the common denominator 12 the numerators are $1, -8, 8, -1$, the table the tests check
against Fornberg's.\cite{fornberg1988} The same computation for Chapra's backward rule for the fourth derivative
on the nodes $-5, \ldots, 0$ gives $-2, 11, -24, 26, -14, 3$ over 1, the row the tests compare with the
published one.\cite[ch.~23]{chapra2015}

\section{Richardson extrapolation}\label{sec:derivations-richardson}

Let $\Dc{h}$ be a base rule with an even expansion, $\Dc{h} = f^{(\dorder)}(x) + a_2 h^2 + a_4 h^4 + \cdots$;
by \cref{sec:derivations-moments} every \code{Central<2>} rule has one. With $h_i = h_0 / 2^i$ and $D_i =
\Dc{h_i}$, define $\Aij{i}{0} = D_i$ and
\begin{equation}
  \Aij{i}{j} = \frac{4^j \Aij{i}{j-1} - \Aij{i-1}{j-1}}{4^j - 1}
             = \Aij{i}{j-1} + \frac{\Aij{i}{j-1} - \Aij{i-1}{j-1}}{4^j - 1}.
  \label{eq:derivations-tableau}
\end{equation}
The claim is that column $j$ has an expansion starting at $h^{2j + 2}$.\cite[Sect.~4.2]{burden2016} By induction
on $j$: suppose $\Aij{i}{j-1} = f^{(\dorder)} + c\, h_i^{2j} + \order(h_i^{2j+2})$ with a coefficient $c$ that does
not depend on $i$, which holds for $j = 1$ with $c = a_2$. The parent one level up has the step $2 h_i$, so
$\Aij{i-1}{j-1} = f^{(\dorder)} + c\, 4^j h_i^{2j} + \order(h_i^{2j+2})$. In the combination $4^j \Aij{i}{j-1} -
\Aij{i-1}{j-1}$ the $h^{2j}$ terms cancel exactly and the constant term is $(4^j - 1) f^{(\dorder)}$; dividing by
$4^j - 1$ leaves $f^{(\dorder)} + \order(h_i^{2j+2})$, again with a leading coefficient independent of $i$. The two
forms in \cref{eq:derivations-tableau} are equal by a line of algebra, and the second is the one the library
computes (\cref{sb:difference-form}).

The a posteriori estimate follows from the same expansion. With $\Aij{i}{j-1} - f^{(\dorder)} = c\, h_i^{2j} +
\cdots$ and $\Aij{i-1}{j-1} - f^{(\dorder)} = 4^j c\, h_i^{2j} + \cdots$,
\[
  \Aij{i}{j} - \Aij{i-1}{j-1}
  = \bigl( \Aij{i}{j} - f^{(\dorder)} \bigr) - \bigl( \Aij{i-1}{j-1} - f^{(\dorder)} \bigr)
  = -4^j c\, h_i^{2j} + \order(h_i^{2j+2}),
\]
so the change of an entry from its parent equals the parent's error to leading order, and the parent's error is
of a lower order in the step than the entry's own, $h_i^{2j}$ against $h_i^{2j+2}$. While truncation dominates,
the change is therefore a conservative estimate of the entry's error. Once rounding dominates it is not: the entries of a
converged column are the limit plus independent rounding of about the size of the rounding floor, and the
smallest of several such differences underestimates the error. This is why \cref{eq:ridders-estimate} takes the
larger of the change and the floor, and why Ridders' method stops once two levels bring no smaller estimate.

\section{The error model and the best step}\label{sec:derivations-model}

The error model of \cref{eq:orders-model} is the sum of the leading truncation term and the rounding term of
\cref{sec:derivations-moments}, with $F = \abs{f^{(\aorder + \dorder)}(x)}$ written for short:
\[
  E(h) = \errc\, h^{\aorder} F + \ampl\, \epsmach \frac{\abs{f}}{h^{\dorder}} .
\]
It is convex in $h > 0$, and its derivative vanishes where
$\aorder \errc\, h^{\aorder - 1} F = \dorder \ampl \epsmach \abs{f} / h^{\dorder + 1}$, that is at
\begin{equation}
  \hbest = \left( \frac{\dorder \ampl\, \epsmach \abs{f}}{\aorder \errc\, F} \right)^{1/(\aorder + \dorder)} .
  \label{eq:derivations-best-step}
\end{equation}
At the minimum the two terms stand in the ratio $\dorder : \aorder$ (truncation to rounding), so
$E(\hbest) = (1 + \aorder / \dorder)\, \ampl \epsmach \abs{f} / {\hbest}^{\dorder}$, which is of order
$\epsmach^{1 - \dorder/(\aorder + \dorder)} = \epsmach^{\aorder/(\aorder + \dorder)}$ times a factor built
from $\abs{f}$, $F$ and the constants. The default step assumes that $f$ varies on the scale $\scale$, meaning
$\abs{f^{(k)}} \approx \abs{f} / \scale^k$, so that $\abs{f} / F \approx \scale^{\aorder + \dorder}$; then
$\hbest \approx \kappa\, \epsmach^{1/(\aorder + \dorder)} \scale$ with the constant
\begin{equation}
  \kappa = \left( \frac{\dorder \ampl}{\aorder \errc} \right)^{1/(\aorder + \dorder)},
  \label{eq:derivations-constant}
\end{equation}
which the library drops: its step is $\epsmach^{1/(\aorder + \dorder)} \scale$ (\cref{eq:orders-step}).

The cost of dropping $\kappa$ is bounded by the model. At a step $h = \hbest / \kappa$ the truncation term is
$\kappa^{-\aorder}$ times its value at $\hbest$ and the rounding term $\kappa^{\dorder}$ times its value, and
since those values are the fractions $\dorder / (\aorder + \dorder)$ and $\aorder / (\aorder + \dorder)$ of
$E(\hbest)$,
\begin{equation}
  \frac{E(\hbest / \kappa)}{E(\hbest)} = \frac{\dorder\, \kappa^{-\aorder} + \aorder\, \kappa^{\dorder}}{\aorder + \dorder} .
  \label{eq:derivations-cost}
\end{equation}
For the first derivative $\kappa$ lies between \num{\valStepConstantCentralTwoFirst} (\code{Central<2>}, where it
is $3^{1/3}$) and 2 (\code{Forward<1>}, $(2 \cdot 2)^{1/2}$), and \cref{eq:derivations-cost} gives at most
\num{\valStepCostCentralTenFirst}, reached by \code{Central<10>} with $\kappa = \num{\valStepConstantCentralTenFirst}$;
for \code{Central<2>} and $f'$ it is \num{\valStepCostCentralTwoFirst}. For the higher orders $\kappa$ is
\num{\valStepConstantCentralTwoSecond} for $f''$ and \num{\valStepConstantCentralTwoFourth} for $f''''$ with
\code{Central<2>}, and the cost by \cref{eq:derivations-cost} is \num{\valStepCostCentralTwoSecond} and
\num{\valStepCostCentralTwoFourth}; for \code{Central<4>} and $f''''$ it is \num{\valStepCostCentralFourFourth}, the
largest of the library's rules. These agree with the design record's figures of 3.5, 11 and 15
(\code{docs/PHASE4_API.md}); the 1.64 that the library's documentation gives for \code{Central<10>} is a rounding
of $\kappa = \num{\valStepConstantCentralTenFirst}$, where the formula gives \num{\valStepCostCentralTenFirst}. They
are model figures for a function that varies on the scale $\scale$ exactly. The design record's prototype
(\code{docs/PHASE4_API.md}) reports the opposite experience on nine test functions that vary faster than $\scale$:
half the default step was better for every order, by a factor 1.2 to 1.8 for \code{Central<2>}, which is why the
library keeps the plain convention of \cref{eq:orders-step} for every order and states its sensitivity instead
(\cref{sb:step-constant}, \cref{sb:fourth-sensitivity}).
